\section{Вторая программа}
\subsection{Один поток}
\subsubsection{Исходный код}
Исходный код второй программы без дополнительных потоков 
представлен в листинге~\ref{lst:prog-2-1}.

\begin{listing}[Исходный код второй программы без дополнительных потоков]{lst:prog-2-1}{C++}
#include <fcntl.h>
#include <unistd.h>

int main()
{
    char c;
    int fl1 = 1, fl2 = 1;

    int fd1 = open("alphabet.txt", O_RDONLY);
    int fd2 = open("alphabet.txt", O_RDONLY);

    while ((fl1 == 1) && (fl2 == 1))
    {
        if ((fl1 = read(fd1, &c, 1)) == 1)
            write(1, &c, 1);
        if (fl1 == 1 && ((fl2 = read(fd2, &c, 1)) == 1))
            write(1, &c, 1);
    }
    return 0;
}

\end{listing}

\subsubsection{Результат работы}
Результат работы второй программы без дополнительных потоков 
представлен на рисунке~\ref{img:2-1.png}.
\jimg{2-1.png}{Результат работы второй программы без дополнительных потоков}

\subsubsection{Анализ полученного результата}
% Системный вызов open создаёт два дескриптора открытого файла в режиме 
% <<только для чтения>> (fd1, fd2). Соответственные дескрипторам структуры 
% file ссылаются на один и тот же inode файла <<alphabet.txt>>. Ссылки на 
% эти структуры помещаются в таблицу открытых процессом файлов под индексами 
% 3 и 4, так как дескрипторы 0,1 и 2 заняты стандартными потоками.

% Системный вызов read для дескриптора fd1 считывает один символ (<<a>>), что 
% приводит к увеличению на 1 значения поля f\_pos в соответствующей структуре 
% file. Считанный символ выводится в стандартный поток вывода.

% Так как дескрипторам fd1 и fd2 соответствуют независимые структуры file, то для 
% fd2 значение поля f\_pos равно 0. Поэтому при вызове read для fd2 будет 
% считан тот же самый первый символ (<<a>>), который также выведется в 
% стандартный поток вывода.

% Поочерёдное чтение и вывод символов будут продолжаться в цикле, пока 
% значения f\_pos в обеих структурах file не достигнут конца файла 
% (размера файла, хранящегося в поле i\_size и равного 26 байтам).

% Таким образом, в стандартный поток вывода каждый символ алфавита 
% будет выведен строго последовательно дважды.

В данном варианте программы процесс выполняет два независимых системных 
вызова open() для одного и того же файла в режиме <<только для чтения>> (O\_RDONLY). 
В результате в таблице открытых файлов процесса резервируются два уникальных 
файловых дескриптора (fd1 и fd2). На уровне ядра операционной системы в 
системной таблице открытых файлов создаются две независимые структуры 
struct file, каждая из которых ссылается на один и тот же inode целевого файла.

Ключевой особенностью данной реализации является то, что каждая из 
созданных структур struct file обладает собственным, независимым от 
другой указателем текущей позиции чтения f\_pos, который при успешном 
открытии файла инициализируется нулем.

В рамках цикла while процесс осуществляет чтение данных, поочередно 
передавая системному вызову read() сначала дескриптор fd1, а затем fd2.

При обращении через fd1 ядро считывает один символ с позиции f\_pos 
первой структуры struct file и инкрементирует её значение на 1.

При последующем обращении через fd2 ядро считывает символ, опираясь 
на f\_pos второй структуры struct file, которая также изначально 
равна нулю, и аналогично увеличивает её на 1.

Поскольку файловые смещения независимы, через оба 
дескриптора процесс читает содержимое файла с самого начала. 
Вывод прочитанных символов в стандартный поток выполняется строго 
последовательно в рамках одного главного потока исполнения. 
Благодаря этому состояние гонки отсутствует, а результат 
программы полностью детерминирован: каждый символ английского 
алфавита из файла выводится на экран ровно дважды подряд 
(формируя последовательность вида <<aabbcc...>>), пока оба 
указателя f\_pos не достигнут конца файла.

\subsection{Один дополнительный поток}
\subsubsection{Исходный код}
Исходный код второй программы с одним дополнительным потоком 
представлен в листинге~\ref{lst:prog-2-2}.

\begin{listing}[Исходный код второй программы с одним дополнительным потоком]{lst:prog-2-2}{C++}
#include <fcntl.h>
#include <pthread.h>
#include <stdio.h>
#include <stdlib.h>
#include <unistd.h>

void *thread_func(void *arg)
{
    char c;
    int fl = 1;
    int fd = *((int *)arg);

    while (fl == 1)
    {
        if ((fl = read(fd, &c, 1)))
            write(1, &c, 1);
    }
    return NULL;
}

int main()
{
    pthread_t t;
    char c;
    int fl = 1;

    int fd1 = open("alphabet.txt", O_RDONLY);
    int fd2 = open("alphabet.txt", O_RDONLY);

    if (pthread_create(&t, NULL, thread_func, &fd1) != 0)
    {
        perror("pthread_create");
        exit(1);
    }

    while (fl == 1)
    {
        if ((fl = read(fd2, &c, 1)) == 1)
            write(1, &c, 1);
    }
    pthread_join(t, NULL);
    return 0;
}
\end{listing}

\subsubsection{Результат работы}
Результат работы второй программы с одним дополнительным потоком 
представлен на рисунке~\ref{img:2-2.png}.
\jimg{2-2.png}{Результат работы второй программы с одним дополнительным потоком}

\subsubsection{Анализ полученного результата}
% Вызов pthread\_create создаёт дополнительный поток, которому в качестве 
% аргумента передаётся указатель на дескриптор fd1. Поскольку на создание 
% потока требуется время, главный поток может успеть прочитать и вывести либо 
% весь алфавит, либо его часть.

% После создания дополнительного потока чтение и вывод символов осуществляются 
% обоими потоками в зависящей от планировщика последовательности.

% Когда главный поток завершит чтение по своему дескриптору, он выйдет из цикла 
% и заблокируется в вызове pthread\_join, ожидая завершения дополнительного потока. 
% Дополнительный поток продолжит работу и выведет оставшуюся часть своих символов.

% Таким образом, каждый символ файла будет напечатан дважды, однако точный 
% порядок их вывода не определён и полностью зависит от того, как планировщик 
% распределяет процессорное время между потоками.

В данной версии программы главный поток первоначально выполняет два независимых 
системных вызова open(), получая два уникальных файловых дескриптора. В ядре 
операционной системы, как и в однопоточном варианте, создаются две независимые 
структуры struct file, каждая из которых ссылается на один и тот же inode и 
инициализируется собственным смещением f\_pos = 0.

Затем главный поток с помощью библиотечной функции pthread\_create() порождает 
один дополнительный поток выполнения, передавая ему указатель на первый 
файловый дескриптор. Главный поток при этом не блокируется и продолжает свое 
исполнение, используя для чтения второй файловый дескриптор. Блокировка главного 
потока функцией pthread\_join() происходит только после завершения его собственного 
цикла чтения.

Оба потока параллельно и независимо друг от друга выполняют системные вызовы 
read(). Поскольку каждый поток оперирует собственной структурой struct file, 
состояние гонки за позицию чтения внутри файла исключено: смещения 
независимы, и каждый из потоков гарантированно прочитает весь файл 
от начала до конца.

Тем не менее, вывод прочитанных символов осуществляется через системный вызов 
write() в стандартный поток вывода, который выступает общим разделяемым ресурсом 
для всего процесса. Поскольку создание нового потока сопряжено с системными 
накладными расходами ядра, главный поток может успеть прочитать и вывести 
некоторую часть данных, в том числе и весь файл до фактического выделения процессорного времени 
дополнительному потоку. 
% В дальнейшем потоки работают конкурентно, и их обращения 
% к терминалы вытесняются и синхронизируются планировщиком операционной системы в 
% произвольном порядке. Из-за этого итоговый вывод содержит два полных набора 
% символов алфавита, однако их последовательность перемешана и строго не детерминирована.

После создания дополнительного потока чтение и вывод символов осуществляются 
обоими потоками в зависящей от планировщика последовательности.

Когда главный поток завершит чтение по своему дескриптору, он выйдет из цикла 
и заблокируется в вызове pthread\_join, ожидая завершения дополнительного потока. 
Дополнительный поток продолжит работу и выведет оставшуюся часть своих символов.

Таким образом, каждый символ файла будет напечатан дважды, однако точный 
порядок их вывода не определён и полностью зависит от того, как планировщик 
распределяет процессорное время между потоками.

\subsection{Два дополнительных потока}
\subsubsection{Исходный код}
Исходный код второй программы с двумя дополнительными потоками 
представлен в листинге~\ref{lst:prog-2-3}.

\begin{listing}[Исходный код второй программы с двумя дополнительными потоками]{lst:prog-2-3}{C++}
#include <fcntl.h>
#include <pthread.h>
#include <stdio.h>
#include <stdlib.h>
#include <unistd.h>

void *thread_func(void *arg)
{
    char c;
    int fl = 1;
    int fd = *((int *)arg);

    while (fl == 1)
    {
        if ((fl = read(fd, &c, 1)))
            write(1, &c, 1);
    }
    return NULL;
}

int main()
{
    pthread_t t1, t2;

    int fd1 = open("alphabet.txt", O_RDONLY);
    int fd2 = open("alphabet.txt", O_RDONLY);

    if (pthread_create(&t1, NULL, thread_func, &fd1) != 0)
    {
        perror("pthread_create");
        exit(1);
    }
    if (pthread_create(&t2, NULL, thread_func, &fd2) != 0)
    {
        perror("pthread_create");
        exit(1);
    }

    pthread_join(t1, NULL);
    pthread_join(t2, NULL);

    return 0;
}
\end{listing}

\subsubsection{Результат работы}
Результат работы второй программы с двумя дополнительными потоками 
представлен на рисунке~\ref{img:2-2.png}.
\jimg{2-3.png}{Результат работы второй программы с двумя дополнительными потоками}

\subsubsection{Анализ полученного результата}
% Два вызова pthread\_create создают два дополнительных потока. Каждому из них 
% передаётся указатель на свой файловый дескриптор (первому --- fd1, 
% второму --- fd2). Главный поток, создав рабочие потоки, сразу же блокируется 
% на вызовах pthread\_join, ожидая их завершения.

% На создание потоков требуется время и их порядок выполнения зависит от 
% планировщика, поэтому порядок вывода символов неизвестен.

% Таким образом, каждый символ файла будет напечатан дважды, однако точный 
% порядок их вывода не определён и полностью зависит от того, как планировщик 
% распределяет процессорное время между дополнительными потоками.

В этой реализации ответственность за открытие файлов возложена исключительно 
на главный поток: до создания дополнительных потоков он дважды выполняет системный 
вызов open(), получая два файловых дескриптора. На каждый вызов в системной таблице 
открытых файлов ядра создается отдельная структура struct file с нулевым 
начальным смещением.

После этого главный поток с помощью вызовов pthread\_create() создает два дополнительных 
рабочих потока, передавая каждому из них по одному уникальному дескриптору. 
Сразу после создания потоков главный поток переходит в заблокированное состояние 
посредством вызовов pthread\_join(), ожидая завершения работы своих потомков.

С точки зрения структур ядра логика работы с файлом остается идентичной предыдущему 
варианту. Два рабочих потока, используя выделенные им дескрипторы, независимо 
и параллельно выполняют системные вызовы read(), считывая данные от начала до 
конца файла. Конфликтов за указатель текущей позиции не происходит, так как 
каждый дескриптор привязан к своей структуре struct file.

Вывод прочитанных символов осуществляется конкурентно через вызовы write() в 
общий стандартный вывод. Поскольку порядок исполнения потоков, их старт и 
выделяемые им кванты процессорного времени полностью контролируются планировщиком 
ОС, точная очередность вывода символов не может быть предсказана. Вследствие 
параллельной работы потоков и отсутствия механизмов синхронизации при обращении к 
stdout, итоговый результат детерминирован лишь по составу данных (каждый символ 
исходного файла будет напечатан ровно дважды), но на экране формируется непредсказуемо 
перемешанная последовательность из двух алфавитов.

\subsection{Схема связи структур}
\jimg{2.pdf}{Схема связи структур в первой программе}